This chapter converts the gap identified in the literature into a focused research plan. It defines the clinical case, research question, hypotheses, objectives, and the limits of the evidence required to answer the question.

\section{Clinical Case for the Study}

Healthcare VLMs may support diagnosis, pathology review, clinical communication, and treatment planning, but a useful human-in-the-loop system must do more than generate fluent answers. It should also provide a signal when an answer should be withheld or referred for expert review. This matters at case level: a high average score can coexist with a small group of confident failures, and those cases may be more important to a clinician than a modest difference between benchmark means. Selective prediction formalises the separation between answering and abstaining \citep{geifman2017selective}, while responsible healthcare AI guidance stresses that errors and their consequences must be evaluated directly \citep{wiens2019donoharm}. Medical image questions make the distinction concrete: sampled answers may be mutually consistent but still be wrong or unrelated to the question, so consistency cannot substitute for validity \citep{carlini2026qasnne}. The intended use of a detector is therefore not to replace clinical judgement, but to provide evidence for deciding which cases require it. This dissertation asks that preclinical evaluation question and does not claim deployment readiness, patient benefit, or clinical safety.

\section{Problem Statement}

Reference metrics describe how closely a greedy answer matches a reference, but they do not show whether the model can identify its own weak outputs. Failure detection uses a different source of evidence: a score derived from sampled answers must rank metric-defined failures above non-failures. The detector score and the failure label must remain separate; otherwise the evaluation would reward a method for reproducing a target that it helped define. Semantic uncertainty can improve over exact string agreement, yet it remains sensitive to the model, task, and data distribution \citep{farquhar2024semanticentropy,ovadia2019uncertainty}. The unresolved problem is whether models that rank highly for answer generation also rank highly for failure detection when the checkpoints, prompts, decoding rules, references, and failure definition are held constant. Without this control, an apparent model difference could come from the evaluation rather than the model. PathVQA, VQA-RAD, and ProstateMM-CHIMERA provide equal but different pathology, radiology, and prostate assessment settings in which to examine that relationship \citep{he2020pathvqa,lau2018vqarad,chimera2025task1}.

\section{Research Question}

The primary research question is: \emph{Are healthcare VLMs with stronger answer-generation performance also better at recognising their unreliable answers in zero-shot medical visual question answering?} For model $m$ and dataset $d$, answer performance is represented by $\mathbf{A}_{m,d}=(\mathrm{BLEU1},\mathrm{BLEU2},\mathrm{ROUGE\mbox{-}L},\mathrm{METEOR})$, while detector $k$ is represented by $\mathbf{R}_{m,d,k}=(\mathrm{AUROC},\mathrm{AUPRC},\mathrm{selective\ performance})$. The analysis asks whether the model ordering produced by $\mathbf{A}$ agrees with that produced by $\mathbf{R}$. Agreement would mean that the strongest answer generator also produces the most informative failure signal in the same setting. Divergence would mean that answer performance cannot stand in for reliability evaluation, while a changing relationship would make the conclusion dataset-dependent. Four supporting questions follow: how the three models rank on answer quality; how well their sampled outputs detect operational failures; whether the two rankings agree within each dataset; and whether the relationship changes across the three medical settings.

\section{Hypotheses}

The main working hypothesis is that answer-generation strength and failure awareness are related but distinct, so their model rankings will not show uniform agreement across all datasets. A second hypothesis is that detector performance depends on the model and dataset because architectures, training, image modalities, question forms, and clinical context differ; uncertainty is known to change under distribution shift \citep{ovadia2019uncertainty}. A third expectation is that an informative detector will improve retained answer quality and reduce the operational failure rate when high-scoring cases are rejected compared with random selection. The first hypothesis is supported by at least one clear ranking divergence, the second by changes in model or detector order across datasets, and the third by improvement over random ranking at matched coverage. Uniform ranking agreement would weaken the first hypothesis, while selective performance close to random would weaken the third. These are descriptive hypotheses about ranking and selective behaviour, not confirmatory tests of patient outcomes.

\section{Aim and Objectives}

The aim is to determine whether stronger zero-shot answer generation can be treated as evidence of stronger operational failure awareness in healthcare VLMs. Six linked objectives answer this aim: (1) evaluate Qwen2.5-VL-3B-Instruct, LLaVA-1.5-7B, and MedGemma-4B-it under a controlled zero-shot protocol; (2) measure greedy answers with BLEU-1, BLEU-2, ROUGE-L, and METEOR; (3) apply one reproducible operational failure rule to all outputs; (4) derive AFD frequency, Semantic AFD, and question-aligned uncertainty from three sampled answers, with a random baseline; (5) evaluate discrimination using AUROC and AUPRC and examine selective performance at fixed coverages; and (6) compare answer-quality and failure-detection rankings within each dataset and across settings. Together, these steps form one evidence chain from controlled generation to answer measurement, failure definition, detector evaluation, and final ranking comparison.

\section{Scope and Definitions}

The study is retrospective, quantitative, comparative, and preclinical. It uses 6,719 PathVQA test items, 451 VQA-RAD test items, and 42 ProstateMM-CHIMERA task records from 14 patients. Each dataset has a domain-specific prompt that is fixed across the three models, and no checkpoint is fine-tuned on the test split. One greedy output is evaluated, while three sampled outputs are used only for detector scores. \emph{Answer quality} means reference agreement under the four stated metrics. An \emph{operational failure} occurs when item-level ROUGE-L is below 20\% and METEOR is below 10\%. \emph{Failure awareness} means that a detector ranks these failures above non-failures; uncertainty is the predictor, not the label. Scores are calculated on their standard normalised scales and displayed as percentages, with differences reported in percentage points.

\section{Expected Contribution}

The contribution is an evaluation framework rather than a new VLM or uncertainty algorithm. It places answer generation and operational failure detection in one controlled comparison without merging them into a single score, compares three differently trained VLMs across three equal medical settings, and links detector discrimination to selective performance. Agreement between the two rankings would suggest that answer strength provides some evidence about failure awareness under the tested conditions; divergence would show that the properties require separate evaluation; and a changing relationship would show that conclusions depend on task and data distribution. Because the failure label is a lexical proxy rather than expert-adjudicated harm, the findings describe metric-defined reliability and cannot establish grounding, calibration for deployment, patient safety, or improved clinical decisions.
